\subsection[Challenge 40: Hydrodynamic modes in Schwinger-Keldysh]{Challenge 40: Hydrodynamic modes in Schwinger-Keldysh}
\textbf{Problem ID:} \texttt{Challenge\_40\_main}\\
\textbf{Primary inferred discipline:} Statistical Physics \& Thermodynamics\\
\textbf{Secondary inferred disciplines:} Condensed Matter Physics\\
\textbf{Python implementation:} \texttt{challenges/40/answer.py}

\subsubsection*{Problem}

\begin{quote}
\color{blue}
\textbf{Reviewer clarification.}
The original symmetry-breaking pattern does not by itself determine a unique
infrared theory: different choices of local multipole densities and Goldstone
fields can produce different spectra. The additions in blue specify the
minimal one-dimensional quadrupole analogue of the dipole-condensate
hydrodynamics discussed by
\href{https://arxiv.org/abs/2303.09573}{Stahl et al.}; the non-uniqueness of
Goldstone realizations for partial multipole breaking is discussed by
\href{https://arxiv.org/abs/2111.08041}{Stahl, Lake, and Nandkishore}.
\end{quote}

\paragraph*{Problem setup}
Effective field theory is a powerful tool used to construct phenomenological
models via symmetries. The method has recently been extended to dissipative
systems via the Schwinger-Keldysh formalism. In this problem we will study the
dissipative effective field theory associated to spontaneous symmetry breaking
of multipolar $U(1)$ symmetries.

{\color{blue}
Consider a $1d$ system with a conserved charge
\[
N=\int_{-\infty}^{\infty}\mathrm{d}x\,\rho(t,x),
\]
a conserved dipole moment
\[
D=\int_{-\infty}^{\infty}\mathrm{d}x\,x\rho(t,x),
\]
and a conserved quadrupole moment
\[
Q=\int_{-\infty}^{\infty}\mathrm{d}x\,x^2\rho(t,x).
\]
Assume that these are the only conserved quantities in the system.
}

Suppose that the quadrupole $Q$ generator is spontaneously broken, and the
charge $N$ and dipole $D$ generators are unbroken.

{\color{blue}
Work to linear order about a homogeneous equilibrium state at fixed
temperature. Write $n=\delta\rho$ for a localized density fluctuation and let
$\mu$ be its conjugate chemical potential. Let $\theta$ be the Goldstone field
associated with the broken quadrupole symmetry. Take $n$ and $\theta$ to be the
only independent hydrodynamic fields.

Let $\chi$ be the charge susceptibility and $\kappa$ the quadrupole superfluid
stiffness. Let $\sigma$ be the Onsager coefficient in the rank-two current
$J^{(2)}$ associated with Goldstone relaxation; assume that the rank-three
current $J^{(3)}$ has no independent dissipative contribution. You may assume
all other EFT coefficients are zero.
}

\paragraph*{Main problem}

Compute the spectrum of hydrodynamic modes $\omega(k)$.

Express your answer in terms of $\chi$, $\kappa$, and $\sigma$.
